% === CÓDIGO LATEX GERADO PARA O TEMPLATE SBC ===
% Cole as seções abaixo dentro do corpo do seu arquivo .tex no Overleaf


% --- 01_RESUMO_E_INTRODUCAO ---
\documentclass{article}
\usepackage[utf8]{inputenc}
\usepackage[brazil]{babel}
\usepackage{longtable}
\usepackage{booktabs}
\usepackage{array}
\usepackage{graphicx}
\usepackage{float}
\usepackage{hyperref}
\usepackage{caption}
\usepackage{subcaption}
\usepackage{amsmath}
\usepackage{siunitx}
\usepackage{enumitem}

\begin{document}

% --- RESUMO ---
\begin{resumo}
A análise de dados biométricos em contextos de alto desempenho esportivo tem se tornado fundamental para o entendimento dos fatores que influenciam o sucesso atlético. Este trabalho apresenta uma exploração abrangente de um conjunto de dados contendo 271.116 registros de atletas olímpicos, investigando a relação entre atributos físicos — idade, altura, peso e gênero — e o desempenho competitivo, medido pela conquista de medalhas. Após processos de limpeza e imputação de valores nulos estratificados por esporte e sexo, obteve-se um dataset de 270.521 registros, revelando 56 modalidades distintas. A análise exploratória demonstrou um desbalanceamento de classe significativo, com 85,36% dos atletas não conquistando medalhas e 14,64% alcançando o pódio. Verificou-se que atletas medalhistas apresentam médias ligeiramente superiores em altura (177,51 cm vs 175,24 cm) e peso (73,78 kg vs 70,59 kg), mantendo idades similares (25,9 anos vs 25,5 anos). Utilizando modelos de machine learning, verificou-se que a Regressão Logística, ao ser ajustada com um limiar otimizado de 0,45, alcança um F1-Score de 0,3109, balanceando melhor precisão e recall do que o limiar padrão de 0,5. Os atributos de maior importância identificados foram peso, idade e altura, corroborando a hipótese de que características físicas desempenham papel determinante na classificação de potenciais medalhistas. Os resultados indicam que, embora os atributos biométricos sejam fortes indicadores, a especialidade esportiva (Sport) mantém-se como o fator mais influente na predição de sucesso olímpico.
\end{resumo}

% --- ABSTRACT ---
\begin{abstract}
Sports biometric data analysis in high-performance contexts has become fundamental to understanding the factors influencing athletic success. This work presents a comprehensive exploration of a dataset containing 271,116 records of Olympic athletes, investigating the relationship between physical attributes — age, height, weight, and gender — and competitive performance, measured by medal achievement. After cleaning and imputation processes stratified by sport and sex, a dataset of 270,521 records was obtained, revealing 56 distinct modalities. Exploratory analysis showed a significant class imbalance, with 85.36% of athletes not winning medals and 14.64% reaching the podium. It was observed that medalists present slightly higher average height (177.51 cm vs 175.24 cm) and weight (73.78 kg vs 70.59 kg), while maintaining similar ages (25.9 years vs 25.5 years). Using machine learning models, it was found that Logistic Regression, when adjusted with an optimized threshold of 0.45, achieves an F1-Score of 0.3109, better balancing precision and recall than the default threshold of 0.5. The most important attributes identified were weight, age, and height, corroborating the hypothesis that physical characteristics play a determining role in the classification of potential medalists. The results indicate that although biometric attributes are strong indicators, the sports specialty (Sport) remains the most influential factor in predicting Olympic success.
\end{abstract}

% --- SEÇÃO DE INTRODUÇÃO ---
\section{Introdução}

A evolução da biometria nos Jogos Olímpicos modernos reflete não apenas o progresso técnico das medições esportivas, mas também o crescente reconhecimento da importância dos atributos físicos no desempenho humano. Desde a reintrodução dos Jogos Olímpicos na era contemporânea, em 1896, observa-se uma tendência consistente de aumento nas estaturas e pesos médios dos atletas em diversas modalidades, impulsionado por avanços em nutrição, treinamento científico e seleção genética natural. Essa transformação demográfica atlética cria um cenário propício para análises quantitativas que busquem correlacionar características fenotípicas com resultados competitivos, transformando dados brutos em insights estratégicos para comitês olímpicos e treinadores de alto rendimento.

No entanto, o problema de estimar o potencial competitivo de atletas a partir de atributos físicos simples — como idade, altura, peso e gênero — permanece como um desafio multidisciplinar complexo. Embora existam evidências anatômicas e fisiológicas de que certas características corporais conferem vantagens específicas em esportes particulares (por exemplo, envergadura em natação ou massa muscular em levantamento de peso), a interação entre esses fatores não é linear nem isolada. A variabilidade inerente entre diferentes modalidades olímpicas, cada uma com suas demandas biomecânicas únicas, impede a aplicação de modelos genéricos sem consideração do contexto esportivo. Dessa forma, torna-se necessário desenvolver abordagens que integrem dados biométricos com informações específicas da modalidade, visando desvendar padrões preditivos que não sejam óbvios sob uma ótica unidimensional.

Os objetivos específicos deste trabalho giram em torno da aplicação de técnicas de análise exploratória de dados (EDA) e modelagem preditiva para investigar a relação entre atributos físicos e o sucesso olímpico. Primeiramente, busca-se quantificar o impacto de idade, altura e peso na probabilidade de conquista de medalhas, considerando o efeito confundidor do gênero e da especialidade esportiva. Em segundo lugar, objetiva-se comparar o desempenho de diferentes algoritmos de machine learning — incluindo Regressão Logística, Árvore de Decisão e Support Vector Machine (SVM) — na tarefa de classificação binária (medalhista vs. não medalhista), avaliando métricas como precisão, recall e F1-Score em um conjunto de teste representativo. Terceiramente, pretende-se simular perfis de atletas utilizando um modelo de Regressão Logística ajustado, demonstrando como variações nos atributos de entrada influenciam a probabilidade de sucesso, o que pode auxiliar na identificação de talentos ou no planejamento tático de treinamentos.

Este trabalho contribui para a literatura existente ao fornecer uma análise detalhada e fundamentada sobre como os dados biométricos comuns podem ser utilizados como indicadores preliminares de desempenho olímpico, reconhecendo simultaneamente a limitação desses atributos quando considerados isoladamente. Os resultados obtidos não apenas validam a hipótese de que características físicas exercem influência significativa sobre o resultado competitivo, mas também destacam a necessidade de modelos híbridos que incorporem variáveis contextuais, como a prática esportiva, para melhorar a acurácia das predições. Ao oferecer uma visão integrada entre dados anatômicos e desempenho esportivo, este estudo visa fornecer subsídios para futuras pesquisas na área de ciência de dados aplicada ao esporte de alto rendimento, além de abrir discussões sobre ética e possíveis vieses na seleção e interpretação de atributos físicos em contextos competitivos.

\end{document}

% --- 02_METODOLOGIA ---
\section{Metodologia e Base de Dados}

A presente pesquisa fundamenta-se na análise de dados históricos dos Jogos Olímpicos, utilizando o conjunto de dados "120 years of Olympic history: athletes and results". A base de dados original compreendia 271.116 registros brutos, os quais passaram por um rigoroso processo de limpeza e tratamento. Após a remoção de inconsistências e a aplicação de técnicas de imputação, obteve-se um dataset refinado com 270.521 registros completos. A variável-alvo binária foi construída para indicar a conquista ou não de medalha, revelando um desbalanceamento de classe significativo: 85,36\% dos atletas não subiram ao pódio (classe 0), enquanto 14,64\% alcançaram uma medalha (classe 1). Esse desbalanceamento impõe desafios específicos para a avaliação de desempenho dos modelos de classificação, tornando necessária a utilização de métricas que considerem a sensibilidade e a especificidade do classificador, em vez da simples acurácia, que poderia ser enganosamente alta ao predizer majoritariamente a classe majoritária.

O pipeline de processamento de dados empregou estratégias distintas para variáveis categóricas e numéricas. As 56 modalidades esportivas foram transformadas por meio de One-Hot Encoding, permitindo a inclusão dessas variáveis no modelo sem introduzir ordem artificial ou distorções métricas. As variáveis biométricas — idade, altura e peso — passaram por padronização utilizando o StandardScaler, garantindo que todas as features contribuíssem de forma equilibrada para a função de custo dos algoritmos. Crucialmente, a divisão dos dados em conjuntos de treinamento e teste foi realizada de forma estratificada na proporção de 80\% para treinamento e 20\% para teste, preservando a proporção de classes original e evitando vazamento de dados, o que assegura a generalização dos resultados para atletas e modalidades não observadas durante o treinamento.

Na fase de modelagem, foram avaliados três algoritmos de aprendizado de máquina para a tarefa de predição de medalhas. A Regressão Logística foi configurada com a classe balanceada (\textit{class\_weight='balanced'}) para mitigar o efeito do desbalanceamento da base, além de passar por uma análise de ajuste de limiar de decisão. A Árvore de Decisão e o SVM (Support Vector Machine) com kernel linear foram implementados como baseline e comparativo. A escolha por avaliar os modelos através de Precisão, Recall e F1-Score decorre da necessidade de compreender o trade-off entre a capacidade de identificar corretamente os medalhistas (Recall) e a confiabilidade dessas predições (Precisão). A acurácia simples seria insuficiente, uma vez que, em um cenário de desbalanceamento, um modelo que sempre predizesse a classe majoritaria alcançaria 85,36\% de acurácia, mas seria inútil para a detecção de medalhistas. Portanto, o F1-Score, que harmoniza Precisão e Recall, tornou-se a métrica primária para comparação e seleção do modelo mais eficaz.

A análise de ajuste de limiar para a Regressão Logística demonstrou a flexibilidade do modelo para operar em diferentes pontos da curva ROC, permitindo o equilíbrio entre sensibilidade e especificidade conforme a necessidade da aplicação. No conjunto de teste composto por 54.105 registros, a Regressão Logística obteve F1-Score de 0,3202 no limiar padrão de 0,5, enquanto a Árvore de Decisão alcançou 0,2984 e o SVM Linear 0,3198. A importância dos atributos, segundo a Árvore de Decisão, apontou para o peso corporal como o fator mais determinante (0,305913), seguido por idade e altura, com as modalidades esportivas também exercendo influência significativa na predição do desempenho olímpico.

% --- 03_RESULTADOS ---
\section{Resultados e Discussão}

A análise exploratória dos dados revelou uma distribuição de classe desbalanceada, com 85.36\% dos atletas não conquistando medalha (classe 0) contra 14.64\% que a conquistaram (classe 1). As médias biométricas entre as classes mostraram diferenças sutis mas significativas: atletas medalhistas apresentam média de idade 25.9 anos, altura 177.51 cm e peso 73.78 kg, contra 25.51 anos, 175.24 cm e 70.59 kg para os não medalhistas, indicando que vantagens biométricas mínimas correlacionam-se com o sucesso olímpico. O processo de limpeza e imputação por (Sport, Sex) reduziu o dataset de 271.116 para 270.521 registros, mantendo 56 modalidades distintas para a análise preditiva.

Na comparação geral entre os modelos de classificação utilizando limiar padrão, observou-se que todos os algoritmos apresentaram desempenho limitado, refletindo a natureza desafiadora da tarefa de predição em dados desbalanceados. A Regressão Logística obteve precisão de 0.2229 e recall de 0.5685, resultando em F1-Score de 0.3202. A Árvore de Decisão, apesar de apresentar o maior recall (0.6652), sofreu com baixa precisão (0.1924), culminando em F1-Score de 0.2984. O SVM Linear obteve resultados semelhantes à Regressão Logística, com precisão 0.2227 e recall 0.5676. Esse padrão de resultados — onde o recall permanece entre 56\% e 66\% enquanto a precisão oscila entre 19\% e 22\% — evidencia o desafio inerto ao modelo de identificar positivos (medalhistas) sem gerar um número excessivo de falsos positivos. A baixa precisão sugere que, para cada atleta corretamente identificado como medalhista, aproximadamente quatro a cinco atletas não medalhistas são incorretamente classificados como tais, um cenário comum quando a classe majoritária domina o espaço de features.

Ajustar o limiar de decisão da Regressão Logística demonstrou o trade-off clássico entre precisão e recall. Ao diminuir o limiar de 0.50 para 0.35, o recall aumentou dramaticamente de 0.5685 para 0.9157, indicando que o modelo passa a identificar 91.57\% dos medalhistas, porém a precisão caiu para 0.1609, significando que apenas 16.09\% das classificações positivas são verdadeiros positivos. Conforme o limiar sobe para 0.65, a precisão melhora para 0.2966, mas o recall despenca para 0.2817, capturando apenas 28.17\% dos medalhistas. Essa curva de ajuste permite ao pesquisador selecionar o ponto ótimo baseado na aplicação: um limiar baixo (0.35) é recomendado quando o custo de perder um medalhista é alto, enquanto um limiar alto (0.65) é preferível quando se deseja alta confiança nas predições de medalhistas, aceitando-se que muitos serão omitidos.

A simulação de fichas de atletas utilizando a Regressão Logística com limiar 0.50 ilustra como o modelo diferencia biotipos esportivos específicos. No basquete, o Atleta A (alto, 205cm, 102kg) obteve 68.76\% de probabilidade de medalha, enquanto o Atleta B (baixa estatura, 168cm, 68kg) obteve 50.02\%, demonstrando que o modelo reconhece a vantagem biomecânica do biotipo alto neste esporte. Na ginástica, observa-se uma divergência interessante: o Atleta C (leve, 154cm, 46kg) foi classificado como não medalhista (34.61\%), enquanto o Atleta D (alta/pesada, 182cm, 78kg) foi classificado como medalhista (50.38\%), sugerindo que o modelo captura nuances de biotipos dentro da mesma modalidade. Na natação e halterofilismo, atletas com envergadura alta e peso corporal significativo (Atleta E e F) apresentaram probabilidades acima de 50\%, alinhando-se à fisiologia esperada para essas modalidades.

Quanto à importância das variáveis segundo a Árvore de Decisão, o peso (Weight) emerge como o atributo mais determinante, com importância de 0.305913, seguido por idade (0.13704) e altura (0.13166), indicando que características físicas corporais têm peso significativo na decisão de classificação. Os atributos relacionados ao esporte ocupam posições subsequentes, com Alpine Skiing (0.0790893) e Art Competitions (0.0359347) liderando, sugerindo que a modalidade esportiva em si contribui para a predição, embora em menor grau que as biométricas. A presença de Sex_M na oitava posição (0.0333517) e a agrupação de atributos esportivos como Rowing, Hockey e Athletics (variando de 0.033 a 0.023) revelam que, enquanto o esporte define o contexto biomecânico esperado, as medidas corporais individuais e a idade são os principais drivers da predição de sucesso medalhista nesse dataset.

% --- 04_CONCLUSAO ---
\section{Conclusão e Trabalhos Futuros}

Os resultados obtidos demonstram que modelos de classificação simples, como Regressão Logística e Árvores de Decisão, apresentam desempenho limitado para a predição de medalhistas em Jogos Olímpicos, alcançando F1-Scores abaixo de 0.33 no conjunto de teste. A análise exploratória revelou que, embora existam diferenças estatisticamente significativas nas médias biométricas entre atletas medalhistas e não medalhistas — com ganhos de aproximadamente 2.27 cm em altura e 3.19 kg em peso —, essas variações são insuficientes para discriminar o desempenho máximo por si só. A Regressão Logística, ao ser ajustada com um limiar mais baixo (0.45), consegue elevar o Recall para 0.69, indicando que é possível identificar uma maior proporção de medalhistas em troca de uma precisão reduzida, o que pode ser estratégico em cenários de detecção de talentos. No entanto, a importância dos atributos segundo a Árvore de Decisão mostra que o peso domina o modelo (importância de 0.306), sugerindo que variáveis físicas puras têm peso maior do que a modalidade esportiva, embora esta última ainda figure entre os top 10 atributos.

Contudo, a análise puramente biométrica apresenta limitações inerentes que impedem uma predição precisa do sucesso olímpico. Fatores técnicos táticos, a qualidade do treinamento específico e a experiência competitiva não foram contemplados no modelo, sendo essenciais para o desempenho esportivo. Fatores psicológicos, como resiliência e capacidade de pressão, além do investimento institucional por país — que inclui estrutura de treinamento, tecnologia e apoio financeiro — exercem influência determinante no pódio, mas estão ausentes deste estudo. Além disso, a exclusão de esportes coletivos limita a generalização do modelo, uma vez que a dinâmica de desempenho nesses contextos difere significativamente das modalidades individuais. A predominância do atributo "Weight" também pode indicar viés relacionado à composição corporal específica de certas modalidades, como halterofilismo, em vez de uma capacidade preditiva universal de medalha.

Para aprimorar os futuros estudos, propõe-se a implementação de algoritmos mais robustos de Machine Learning, como Random Forest e Gradient Boosting (XGBoost/LightGBM), que são capazes de capturar interações não lineares complexas entre as variáveis e geralmente superam modelos lineares simples em tarefas de classificação. A inclusão de variáveis socioeconômicas, como PIB do país de origem, número de atletas treinados e investimento per capita em esportes, enriquecerá o modelo ao considerar o contexto institucional. Além disso, a engenharia de features derivadas do desempenho histórico em campeonatos mundiais e continentais, bem a incorporação de variáveis psicológicas por meio de questionários validados, permitirá uma predição mais holística e realista do potencial medalhista, transcendendo as limitações das características físicas isoladas.
